\documentclass[letterpaper,10pt,conference]{ieeeconf}

\IEEEoverridecommandlockouts
\overrideIEEEmargins

\usepackage[T1]{fontenc}
\usepackage{amsmath,amssymb}
\usepackage[hidelinks]{hyperref}

% Keep the ieeeconf.cls supplied with this file in the same directory.
% This draft combines the two supplied texts; references and results can
% be added later. Author and affiliation information remains to be added.
\title{\LARGE\bfseries
Hierarchical Learning and Control for Synchronized Arrival
under a Shared Deadline}
\author{}
\date{}

\begin{document}
\maketitle
\thispagestyle{empty}
\pagestyle{empty}

\begin{abstract}
We present a hierarchical learning and control framework for synchronized
arrival of mobile robots at arbitrary preassigned goals under a shared
deadline. A shared policy, trained using Multi-Agent Proximal Policy
Optimization (MAPPO) with centralized training and decentralized execution,
generates local waypoints and preferred speeds from each vehicle's own state,
goal, timing information, and locally sensed neighbours. Pure-pursuit rollout
converts these decisions into timed reference trajectories, which independent
nonlinear model predictive control (NMPC) controllers track while accounting
for actuator limits and predicted inter-vehicle separation. The controller
operates inside the training loop, allowing the policy to learn coordination
under the constraints and interventions encountered during execution. A
masked attention encoder aggregates the nearest neighbours into a fixed-size
representation, allowing the policy architecture to accommodate varying
neighbourhood sizes. Scalability and robustness are assessed empirically.
Online decisions require neither a centralized planner nor exchanged
neighbour goals or future trajectories, given the supplied common deadline.
We assess the contributions of learning and constrained control through
matched comparisons of pure-pursuit execution, MPC added after policy
training, and training with MPC in the loop, alongside a non-learning deadline
controller. These experiments characterize the trade-offs between synchronized
arrival, collision avoidance, and completion failures.
\end{abstract}

\section{Introduction}
\label{sec:introduction}

The contribution is learning decentralized motion references that coordinate
both route selection and speed so that mobile robots reach different assigned
goals at a common deadline, while accounting for the behaviour of a
constrained low-level controller.

The architecture uses a hierarchical combination of reinforcement learning
(RL) and nonlinear model predictive control (NMPC) to address the problem of
simultaneous arrival, with robustness and scalability evaluated experimentally.

% Expand the introduction and add related work with verified citations later.

\section{Background and Preliminaries}
\label{sec:background}

\subsection{Decentralized Partially Observable Markov Decision Process}

Multi-agent reinforcement learning (MARL) is based on reinforcement learning,
in which an agent interacts with its environment and learns a policy that
aims to maximize the expected cumulative reward.

The problem is modelled as a decentralized partially observable Markov
decision process (Dec-POMDP). In this work, with homogeneous agents and
deterministic local observations, we write the Dec-POMDP as
\begin{equation}
    \mathcal{M}
    = (\mathcal{S},\mathcal{A},\mathcal{O},R,P,\Omega,n,\gamma),
\end{equation}
where $\mathcal{S}$ is the global state space, $\mathcal{A}$ is the common
action space of each agent, $\mathcal{O}$ is the local observation space,
and $n$ is the number of agents. The local observation of agent $i$ at time
step $t$ is
\begin{equation}
    o_{i,t} = \Omega(s_t,i).
\end{equation}
The joint action is
\begin{equation}
    \mathbf{a}_t = (a_{1,t},\ldots,a_{n,t}) \in \mathcal{A}^{n}.
\end{equation}
The transition model $P(s_{t+1}\mid s_t,\mathbf{a}_t)$ describes the
probability of transitioning from $s_t$ to $s_{t+1}$ under the joint action.
The shared reward is $r_t=R(s_t,\mathbf{a}_t)$, and $\gamma\in[0,1)$ is
the discount factor.

Agents use a shared policy $\pi_\theta(a_{i,t}\mid o_{i,t})$, parameterized
by $\theta$, to produce an action from their own local observation. They
jointly optimize the expected discounted cumulative reward:
\begin{equation}
    J_{\mathrm{RL}}(\theta)
    = \mathbb{E}_{\pi_\theta,P}\!\left[
        \sum_{t=0}^{H_{\mathrm{ep}}-1}\gamma^t r_t
      \right],
    \label{eq:rl-objective}
\end{equation}
where $H_{\mathrm{ep}}$ is the number of steps in an episode.

\subsection{Centralized Training with Decentralized Execution}

MARL algorithms can be categorized according to the information available
during training and execution. During training, an algorithm may be
restricted to the local information observed by each agent
(decentralized training), or it may use information about all agents in
the multi-agent system (centralized training).

Centralized training with decentralized execution (CTDE) combines these
information settings: training can use global information, while each agent
uses only its permitted local information to make decisions during execution.

Multi-Agent Proximal Policy Optimization (MAPPO) is an on-policy deep
reinforcement learning algorithm that adopts the CTDE paradigm. In our
implementation, a shared actor $\pi_\theta(a_{i,t}\mid o_{i,t})$ receives
only agent $i$'s local observation and outputs that agent's action. A
centralized critic with parameters $\phi$ receives the global state $s_t$,
which is not supplied to the actor during execution. The critic is used
during training to estimate values and compute advantages using generalized
advantage estimation (GAE). We write its output as $V_\phi(s_t,i)$ because
the implementation produces one value estimate per agent from the global
state.

\subsection{Rollout Collection and Policy Probability Ratio}

Initially, the algorithm collects a sequence of experiences called a
rollout. It stores local observations, global states, sampled latent actions,
old action log-probabilities, critic values, rewards, next-state values,
termination and rollout-boundary flags, and active-agent masks. The latent
Gaussian actions are transformed using $\tanh$ to obtain bounded policy
actions. The policy log-probabilities account for this transformation.

Proximal Policy Optimization (PPO) needs to measure how much the new policy
differs from the policy that generated the rollout. It therefore uses the
probability ratio
\begin{equation}
    \rho_{i,t}(\theta)
    = \frac{\pi_\theta(a_{i,t}\mid o_{i,t})}
            {\pi_{\theta_{\mathrm{old}}}(a_{i,t}\mid o_{i,t})},
    \label{eq:probability-ratio}
\end{equation}
or, equivalently,
\begin{equation}
\begin{aligned}
    \rho_{i,t}(\theta)
    = \exp\!\bigl(&\log\pi_\theta(a_{i,t}\mid o_{i,t})\\
                 &-\log\pi_{\theta_{\mathrm{old}}}(a_{i,t}\mid o_{i,t})\bigr),
\end{aligned}
\end{equation}
where $\theta_{\mathrm{old}}$ denotes the parameters used to collect the
rollout. Both probabilities are evaluated for the same collected action
and observation.

\subsection{Returns and Generalized Advantage Estimation}

The discounted return from time step $t$ is
\begin{equation}
\begin{aligned}
    G_t
    &= \sum_{k=t}^{H_{\mathrm{ep}}-1}\gamma^{k-t}r_k\\
    &= r_t + \gamma r_{t+1} + \gamma^2 r_{t+2} + \cdots.
\end{aligned}
    \label{eq:discounted-return}
\end{equation}
For the collected rollout, GAE estimates how much better an action was
than the critic's value prediction. First, we compute the temporal-difference
error:
\begin{equation}
\begin{aligned}
    \delta_{i,t}
    ={}& r_t + \gamma(1-d_t)V_{\phi_{\mathrm{old}}}(s_{t+1},i)\\
       &-V_{\phi_{\mathrm{old}}}(s_t,i),
\end{aligned}
    \label{eq:td-error}
\end{equation}
where $d_t=1$ if the task terminates after the transition, and $d_t=0$
otherwise. Thus, a true terminal transition does not bootstrap from a
next-state value. The subscript $\phi_{\mathrm{old}}$ indicates that
these are the fixed critic estimates collected before the update.

The GAE advantages are computed backwards through the rollout:
\begin{equation}
\begin{aligned}
    \widehat{A}_{i,t}
    &= \delta_{i,t}
      + \gamma\lambda(1-d_t)(1-b_t)\widehat{A}_{i,t+1},\\
    \widehat{A}_{i,H}&=0,
\end{aligned}
    \label{eq:gae}
\end{equation}
where $H$ is the rollout length, $\lambda\in[0,1]$ controls the GAE
weighting, and $b_t=1$ marks an artificial rollout or curriculum boundary.
Such a boundary stops the advantage recursion but retains the next-state
bootstrap in~\eqref{eq:td-error} when the task has not terminated.

\subsection{Critic Loss and PPO Clipped Objective}

The critic's regression target is the GAE-based return estimate
\begin{equation}
    \widehat{G}_{i,t}
    = \widehat{A}_{i,t} + V_{\phi_{\mathrm{old}}}(s_t,i).
    \label{eq:value-target}
\end{equation}
The critic is trained by minimizing the value loss
\begin{equation}
    \mathcal{L}_{V}(\phi)
    = \frac{1}{2}\,
      \mathbb{E}_{(i,t)\sim\mathcal{B}}\!\left[
        \bigl(V_\phi(s_t,i)-\widehat{G}_{i,t}\bigr)^2
      \right],
    \label{eq:critic-loss}
\end{equation}
where $\mathcal{B}$ denotes the collected agent--time samples. The return
targets are held fixed during the optimization update.

The actor uses the advantages to form the PPO clipped objective. Let
$\widetilde{A}_{i,t}$ denote the advantages normalized over active decision
samples when their variance is nonzero, and define
\begin{equation}
    \overline{\rho}_{i,t}(\theta)
    = \operatorname{clip}\bigl(\rho_{i,t}(\theta),1-\epsilon,1+\epsilon\bigr).
\end{equation}
Then
\begin{equation}
\begin{aligned}
    \mathcal{L}_{\mathrm{clip}}(\theta)
    = \mathbb{E}_{(i,t)\sim\mathcal{B}_{\mathrm{active}}}\!\Bigl[
        \min\bigl(&\rho_{i,t}(\theta)\widetilde{A}_{i,t},\\
                 &\overline{\rho}_{i,t}(\theta)\widetilde{A}_{i,t}\bigr)
      \Bigr],
\end{aligned}
    \label{eq:ppo-clipping}
\end{equation}
where $\epsilon>0$ is the clipping parameter and
$\mathcal{B}_{\mathrm{active}}$ contains samples from agents that were
active when their actions were selected. The actor maximizes this objective.
Clipping limits the incentive for large policy changes in the surrogate
objective; it does not impose a hard bound on every probability ratio.

The implementation combines the actor objective, critic loss, and an entropy
bonus into the loss that is minimized:
\begin{equation}
\begin{aligned}
    \mathcal{L}(\theta,\phi)
    ={}&-\mathcal{L}_{\mathrm{clip}}(\theta)
        +c_V\mathcal{L}_{V}(\phi)\\
       &-c_H\,\mathbb{E}_{(i,t)\sim\mathcal{B}_{\mathrm{active}}}\!\left[
          \mathcal{H}\bigl(\pi_\theta(\cdot\mid o_{i,t})\bigr)
        \right],
\end{aligned}
    \label{eq:combined-loss}
\end{equation}
where $c_V$ and $c_H$ weight the value loss and policy entropy, respectively.
The entropy term encourages exploration. The actor and critic are updated
using minibatches from the rollout, while the stored old log-probabilities,
advantages, and return targets remain fixed during these updates.

\section{Method: Hierarchical Control Pipeline}
\label{sec:method}

The complete control pipeline is
\begin{equation}
\begin{gathered}
\text{Observation}\rightarrow\text{MAPPO}\rightarrow\text{Pure pursuit}\\
\rightarrow\text{NMPC}\rightarrow\text{Vehicle}.
\end{gathered}
\end{equation}
The three main frequencies are $5~\mathrm{Hz}$ for the MAPPO actor,
$20~\mathrm{Hz}$ for NMPC, and $100~\mathrm{Hz}$ for the plant simulation.
The low-level controller operates inside the training loop so that the actor
learns under the controller's constraints and interventions.

In this section, $\mathbf{x}_k$ denotes one vehicle's state, whereas $s_t$
in Section~\ref{sec:background} denotes the global MARL state. The index $k$
denotes a local prediction step within a controller horizon.

\subsection{MAPPO Action and Waypoint Generation}

Every $0.2~\mathrm{s}$, the shared MAPPO actor receives the vehicle's current
local observation, including its state, goal, deadline, schedule information,
and nearby vehicles. The bounded action is
\begin{equation}
    a_t^{\mathrm{RL}}
    = \begin{bmatrix}
        a_{0,t}^{\mathrm{RL}} & a_{1,t}^{\mathrm{RL}} & a_{2,t}^{\mathrm{RL}}
      \end{bmatrix}^{\!\top},
    \qquad a_{q,t}^{\mathrm{RL}}\in[-1,1].
\end{equation}
The action components are decoded as
\begin{align}
    d_f &= 0.25 + \frac{a_{0,t}^{\mathrm{RL}}+1}{2}(1.75),\\
    d_l &= 1.2a_{1,t}^{\mathrm{RL}},\\
    v_{\mathrm{pref}} &= \frac{a_{2,t}^{\mathrm{RL}}+1}{2}(2.0).
\end{align}
Here, $d_f\in[0.25,2.0]~\mathrm{m}$ is the forward waypoint distance,
$d_l\in[-1.2,1.2]~\mathrm{m}$ is the lateral waypoint offset, and
$v_{\mathrm{pref}}\in[0,2.0]~\mathrm{m/s}$ is the preferred speed.
Thus, the actor specifies a local waypoint and a preferred speed.

\subsection{Pure-Pursuit Reference Generation}

The vehicle state is
\begin{equation}
    \mathbf{x}_k
    = \begin{bmatrix}
        p_{x,k} & p_{y,k} & \psi_k & v_k & \delta_k
      \end{bmatrix}^{\!\top},
\end{equation}
where $p_{x,k}$ and $p_{y,k}$ denote position, $\psi_k$ is heading,
$v_k$ is speed, and $\delta_k$ is the steering angle.
The reference rollout starts from the current measured state:
\begin{equation}
    \mathbf{x}^{\mathrm{ref}}_0 = \mathbf{x}_{\mathrm{measured}}.
\end{equation}

Using the MAPPO waypoint, pure pursuit computes a target steering angle.
Before applying steering limits, its geometric expression is
\begin{equation}
    \delta_{\mathrm{target},k}
    = \tan^{-1}\!\left(\frac{2L y_{\mathrm{local},k}}{d_k^2}\right),
\end{equation}
where $L$ is the wheelbase, $d_k$ is the distance from the predicted vehicle
position to the current pursuit target, and $y_{\mathrm{local},k}$ is the
target's lateral coordinate in the predicted vehicle frame. Initially,
when pursuing the decoded waypoint, $y_{\mathrm{local},0}=d_l$. These
relative quantities are recomputed as the reference rollout advances.

The longitudinal acceleration is generated by speed feedback:
\begin{equation}
    a_k^{\mathrm{PP}} = 4\bigl(v_{\mathrm{target},k}-v_k^{\mathrm{ref}}\bigr),
\end{equation}
where $v_{\mathrm{target},k}$ is the preferred speed after the reference
generator's speed caps. The steering-rate command moves the steering angle
towards $\delta_{\mathrm{target},k}$, subject to the control limits.

Pure pursuit propagates the bicycle model for $24$ steps:
\begin{equation}
    \mathbf{x}_{k+1}^{\mathrm{ref}}
    = f(\mathbf{x}_k^{\mathrm{ref}},u_k^{\mathrm{PP}}),
    \qquad k=0,\ldots,23,
\end{equation}
with $\Delta t=0.05~\mathrm{s}$. Therefore, it generates a
$1.2~\mathrm{s}$ timed reference:
\begin{equation}
    X^{\mathrm{ref}}
    = \bigl[\mathbf{x}_0^{\mathrm{ref}},
             \mathbf{x}_1^{\mathrm{ref}},\ldots,
             \mathbf{x}_{24}^{\mathrm{ref}}\bigr].
\end{equation}
The rollout also produces the corresponding nominal controls
$[u_0^{\mathrm{PP}},\ldots,u_{23}^{\mathrm{PP}}]$.
Each reference state specifies where the vehicle should approximately be
at a particular future time, rather than only a point along an untimed path.

\subsection{NMPC Control Input and Prediction}

The NMPC control input is
\begin{equation}
    u_k = \begin{bmatrix}a_k & \dot{\delta}_k\end{bmatrix}^{\!\top},
\end{equation}
where $a_k$ is longitudinal acceleration and $\dot{\delta}_k$ is steering
rate. This is the low-level command optimized by NMPC and is distinct from
the MAPPO waypoint-and-speed action.

The NMPC horizon contains $N=16$ control steps. Since
$\Delta t=0.05~\mathrm{s}$, its duration is
\begin{equation}
    T_H=N\Delta t=16(0.05)=0.8~\mathrm{s}.
\end{equation}
NMPC optimizes the control sequence
\begin{equation}
    U=[u_0,u_1,\ldots,u_{N-1}]\in\mathbb{R}^{2\times N}.
\end{equation}
There are $32$ scalar decision variables because each of the $16$ controls
contains two components.

For each candidate $U$, prediction starts from the current measured state
and propagates the nonlinear bicycle model:
\begin{align}
    \mathbf{x}_0 &= \mathbf{x}_{\mathrm{measured}},\\
    \mathbf{x}_{k+1} &= f(\mathbf{x}_k,u_k),
    \qquad k=0,\ldots,N-1.
\end{align}
The resulting trajectory is
\begin{equation}
    X(U)=[\mathbf{x}_0,\mathbf{x}_1,\ldots,\mathbf{x}_N].
\end{equation}
The predicted states are not independent optimization variables. They are
generated from the candidate controls through the bicycle model. This is
a single-shooting formulation.

\subsection{NMPC Objective Function}

Let $p_k=[p_{x,k},p_{y,k}]^\top$. NMPC compares its predicted states with
the pure-pursuit reference states at the corresponding future times.
Define the errors
\begin{align}
    e_{p,k} &= p_k-p_k^{\mathrm{ref}},\\
    e_{v,k} &= v_k-v_k^{\mathrm{ref}},\\
    e_{\delta,k} &= \delta_k-\delta_k^{\mathrm{ref}},\\
    e_{\psi,k} &= \operatorname{atan2}\bigl(
       \sin(\psi_k-\psi_k^{\mathrm{ref}}),
       \cos(\psi_k-\psi_k^{\mathrm{ref}})\bigr).
\end{align}
The wrapped heading error accounts for the equivalence of angles differing
by $2\pi$. All reference quantities here belong to the time-aligned window
used by the current NMPC solve.

The objective function is
\begin{equation}
\begin{aligned}
    J_{\mathrm{MPC}}(U)=\sum_{k=0}^{N-1}\Bigl(
        &8\|e_{p,k+1}\|^2+3e_{v,k+1}^{2}\\
        &+0.5e_{\psi,k+1}^{2}+0.1e_{\delta,k+1}^{2}\\
        &+0.015\|u_k\|^2\\
        &+0.04\|u_k-u_{k-1}\|^2\Bigr)\\
        &+J_{\mathrm{terminal}},
\end{aligned}
    \label{eq:mpc-cost}
\end{equation}
where $u_{-1}$ is the control command applied during the previous NMPC
step. The terminal cost gives additional importance to the final predicted
position and speed:
\begin{equation}
    J_{\mathrm{terminal}}=16\|e_{p,N}\|^2+6e_{v,N}^{2}.
    \label{eq:terminal-cost}
\end{equation}

\subsection{Why NMPC Can Differ from Pure Pursuit}

Pure pursuit generates a nominal trajectory under the bicycle model and
its control limits. If NMPC used only a tracking-error objective, started
from the same state, and used the same model, the nominal pure-pursuit
controls could reproduce the reference exactly, provided that they also
satisfied all NMPC constraints.

However, the objective in~\eqref{eq:mpc-cost} also includes
\begin{equation}
    0.015\|u_k\|^2+0.04\|u_k-u_{k-1}\|^2.
\end{equation}
The first term penalizes large acceleration and steering-rate commands.
The second penalizes large changes between consecutive controls. NMPC
may therefore accept a small tracking error in exchange for a smoother,
lower-effort control sequence. Separation constraints can also require a
departure from the nominal reference. Consequently, the NMPC command need
not equal the pure-pursuit command even when both use the same bicycle model.

Pure pursuit computes how to steer towards the current target. NMPC instead
optimizes a sequence of controls over the next $0.8~\mathrm{s}$, balancing
tracking, effort, smoothness, and constraints.

\subsection{Collision-Avoidance Constraints}

NMPC also considers nearby vehicles. For vehicle $i$ and a sensed neighbour
$j$, the separation constraint can be written schematically as
\begin{equation}
    d_{ij,k}\geq d_{\mathrm{safe}},
\end{equation}
where $d_{ij,k}$ denotes predicted separation and $d_{\mathrm{safe}}$ is the
required safety distance. Let $\mathcal{N}_i$ denote the sensed neighbours
considered by vehicle $i$. The constrained problem has the form
\begin{equation}
\begin{aligned}
    U^*\in\operatorname*{arg\,min}_{U}\quad
        &J_{\mathrm{MPC}}(U)\\
    \text{subject to}\quad
        &\mathbf{x}_0=\mathbf{x}_{\mathrm{measured}},\\
        &\mathbf{x}_{k+1}=f(\mathbf{x}_k,u_k),\\
        &u_{\min}\leq u_k\leq u_{\max},\\
        &d_{ij,k+1}\geq d_{\mathrm{safe}},\quad j\in\mathcal{N}_i,
\end{aligned}
    \label{eq:mpc-problem}
\end{equation}
for $k=0,\ldots,N-1$. In the single-shooting implementation, the dynamics
are enforced by the rollout itself. The separation constraints apply to
predicted states at the sampled horizon times.
% Specify the exact body geometry, neighbour prediction, state bounds,
% feasibility tolerance, and failure handling from the final implementation.

\subsection{Optimization Process}

The optimizer starts from an initial candidate control sequence $U^{(0)}$.
For this candidate, the prediction model and objective produce
\begin{equation}
    U^{(0)}\rightarrow X(U^{(0)})\rightarrow J_{\mathrm{MPC}}(U^{(0)}).
\end{equation}
The optimizer then modifies the controls through successive candidates
$U^{(1)},U^{(2)},\ldots$, seeking to reduce the objective while satisfying
the constraints. Thus, each candidate is evaluated by propagating the bicycle
model, comparing the predicted trajectory with the timed reference, and
checking the constraints.

The resulting optimized plan is
\begin{equation}
    U^*=[u_0^*,u_1^*,\ldots,u_{N-1}^*].
\end{equation}
Since the problem is nonlinear, the numerical solver generally seeks a local
solution; the notation $U^*$ does not imply a certified global optimum.

\subsection{Receding-Horizon Control}

Although NMPC calculates $16$ future controls, only the first is applied:
\begin{equation}
    u_{\mathrm{applied}}=u_0^*
    =\begin{bmatrix}a_0^* & \dot{\delta}_0^*\end{bmatrix}^{\!\top}.
\end{equation}
The remaining controls allow the optimizer to evaluate the future
consequences of the first command. After $0.05~\mathrm{s}$, the vehicle state
is measured again, the prediction is initialized with the new measurement,
and the optimization is solved again.

This repeated sequence of measurement, prediction, optimization, and
application of the first command is called receding-horizon control.

\subsection{Timing of the Complete Pipeline}

At $t=0$, MAPPO generates $[d_f,d_l,v_{\mathrm{pref}}]$, pure pursuit
generates a $1.2~\mathrm{s}$ reference, and NMPC solves its
$0.8~\mathrm{s}$ optimization problem. NMPC solves again at
$t=0.05$, $0.10$, and $0.15~\mathrm{s}$ using the latest measured state
and the corresponding time-aligned reference window. At $t=0.20~\mathrm{s}$,
MAPPO generates a new action and pure pursuit generates a new reference.

Thus, each policy interval contains four NMPC updates. Each NMPC interval
contains five plant simulation steps of $0.01~\mathrm{s}$. The new measured
vehicle state supplies feedback to the next controller update.

\subsection{Division of Responsibilities}

MAPPO decides where to move and how fast, using the vehicle's permitted local
information. Pure pursuit converts that decision into a timed nominal
trajectory under the bicycle model. NMPC tracks the trajectory while
balancing smoothness and control effort and enforcing its constraints on
the predicted motion. The plant evolves under the applied command and
provides the next measured state.

The hierarchy therefore maps a coordination decision to a timed reference
trajectory and then to a physical control command.

% Add experimental setup, results, discussion, conclusion, and bibliography
% when those parts of the paper are ready.

\end{document}
